\chapter{Introduction}
\label{ch:intro}
\epigraph{It is a laborious madness and an impoverishing one, the madness of composing vast books---setting out in five hundred pages an idea that can be perfectly related orally in five minutes}{Jorge Luis Borges, Prologue to The Garden of Forking Paths}

\lettrine{T}{his thesis} asks how an optimizing compiler can be verified.
A production compiler is a large program: in LLVM~\citep{aho1985compilers,lattner2004llvm}, the peephole rewriter alone is larger than the entire vectorizer!
The tradition of verified compilation~\citep{leroy2016compcert,kumar2014cakeml} teaches that verifying even a modest compiler end to end is the work of many person-years,
so there is no hope for one person to verify all of LLVM.
I therefore focus on the peephole rewriter~\citep{mckeeman1965peephole},
whose unit of work is a single, local transformation on a finite window of instructions.
Each transformation is small enough to state precisely and prove correct,
which has long made the rewriter the natural beachhead for verification,
as witnessed by the Alive line of work on LLVM's instruction combiner~\citep{lopes2015provably,lopes2021alive2}.
But what the individual rewrite grants in simplicity, the rewriter takes back in volume:
a production compiler carries thousands of rewrites, and hence thousands of proof obligations.

Can we prove them by hand? We can, and for a while I did,
when I set out to verify an SSA-based peephole rewriter for LLVM's bitvector optimizations.
Together with collaborators, I formalized a core calculus for SSA-based compiler IRs~\citep{appel1998ssa,rastello2022ssa},
generic over the user's IR, and proved a generic peephole rewriter correct,
with case studies ranging from fully homomorphic encryption to LLVM's bitvector rewrites (\cref{ch:lean-mlir}).
The last of these is where the real difficulty began.
The framework gave us the structural half of the problem for free:
manipulating SSA, threading contexts through rewrites, splicing programs back together.
The arithmetic half it could not give us.
Each rewrite bottoms out in a bitvector identity;
the identities are subtle, hand proofs of them are slow and error-prone,
and at the scale of the thousands of rewrites inside a production compiler, proving them one by one is out of the question.
It is this tedium that drives the demand for push-button automation.

What style of decision procedure should we bet on?
I bet on bitblasting, almost by accident, since I had been helping to implement \bvdecide{}~\citep{leanbv},
Lean's fixed-width bitvector tactic.
\bvdecide{} carries a bitvector goal all the way down to a SAT instance,
dispatches it to a state-of-the-art SAT solver,
and has the resulting certificate rechecked by Lean's kernel, so that no bug in the solver can mint a theorem.
The design deliberately mirrors Bitwuzla~\citep{niemetz2023bitwuzla},
a state-of-the-art SMT solver for fixed-width bitvectors, floating point, and arrays;
our verified infrastructure thus shares its bones with one of the fastest solvers of today,
and it is the bedrock on which the rest of this thesis builds.

However, the proof obligations a compiler hands us are not bitblastable as they stand.
An LLVM rewrite must be correct at \emph{every} register width---at \texttt{i8}, at \texttt{i32}, at \texttt{i1729}---whereas a bitblaster wants one fixed width~\citep{niemetz2019towards,bit-precise-parametric-bv}.
The standard remedy, used for example by Alive, is to enumerate the widths and bitblast each one separately.
We show that this is unnecessary.
Predicates that must hold at all widths reduce to reachability in a finite automaton,
which we decide by $k$-induction and language emptiness (\cref{ch:mono-width}).
Next, by an encoding trick that writes widths as bitmasks,
any formula over $n$ symbolic widths reduces to an equisatisfiable formula at a single width (\cref{ch:multi-width}),
trading an exponential number of solver queries for a \emph{single} query
and carving out a new sound and complete fragment of the multi-width theory.
The reach of bitblasting, it turns out, extends all the way to the parametric bitvector theory.

The last, and by far the most intricate, piece missing from a fully verified Bitwuzla was floating point~\citep{ieeefp,fphandbook}.
The venerable SymFPU~\citep{symfpu} had shown, in a sequence of papers,
that bitblasting decisively outcompetes the alternatives for floating-point SMT solving.
Following its lead, we mechanized the SMT-LIB \qffp{} theory~\citep{smtlibfpa}
and built verified bitblasting circuits for almost all of its operations,
proved correct at every combination of exponent and significand width (\cref{ch:floating-point}).
The width-genericity is no flourish:
machine learning is pushing small, exotic formats such as E4M3 and E5M2 into production hardware~\citep{fp8fordl},
and these are precisely the formats on which existing solvers are least tested and most fragile.

Why should bitblasting, of all approaches, be the one that pays off?
I read it as a flavour of the bitter lesson~\citep{sutton2019bitter},
the observation from AI that search and learning are the only two paradigms that keep rewarding the compute we feed them.
Automated reasoning has its own incarnation of search in SAT solving,
which now routinely settles instances with hundreds of millions of variables~\citep{biere2024cadical,kissatSATComp2024},
and bitblasting~\citep{tseitin1983complexity} is nothing but the craft of encoding richer theories into this engine,
the same bet that powers the state-of-the-art SMT solvers~\citep{niemetz2023bitwuzla,barbosa2022cvc5}.
The experiments in this thesis bear the lesson out:
our width-generic solvers prove more all-width problems than the unverified state of the art,
and our floating-point bitblaster, while slower than Bitwuzla,
offers in exchange guarantees that reach down to Lean's kernel.

There is a second bet running through this thesis: mechanization.
The solvers above are not merely implemented; they are written in Lean~\citep{demoura2021lean} and proved correct in its proof language.
Is this worth the effort? I believe it will only grow in importance.
We live in times (as of August 2026) in which large language models~\citep{brown2020language,touvron2023llama}
produce impressive volumes of code that cannot be trusted to get the details right,
and which regularly fulfil the letter of a specification while betraying its spirit;
a machine-checked specification, with machine-checked automation to discharge it, is the natural antidote.
The lesson of this thesis is that foundational proof and push-button automation need not be at odds,
and I hope the reductions developed here lay the groundwork for verification efforts of far greater ambition.

\section{Contributions}

\noindent First, we contribute \textbf{A verified framework for peephole rewriting in SSA-based IRs} (\cref{ch:lean-mlir}).
We formalize a core calculus for SSA-based IRs~\citep{appel1998ssa,rastello2022ssa},
generic over the user's IR and covering MLIR-style regions~\citep{lattner2020mlir}, mechanize it in
Lean with an MLIR-syntax frontend, and prove correct a generic peephole
rewriter as well as dead-code and common-subexpression elimination.
\publishedat{Published @ ITP~2024~\citep{leanmlir}}

\noindent Second, we develop \textbf{Certified decision procedures for width-independent bitvector predicates} (\cref{ch:mono-width}).
We mechanize key model-checking algorithms, such as $k$-induction, automata reachability, emptiness, and
minimization, and use them to build scalable, certified tactics for
bitvector predicates that hold at \emph{all} widths, plus a specialized
solver for mixed Boolean-arithmetic~\citep{liu2021mba,reichenwallner2022efficient}.
\publishedat{Published @ OOPSLA~2025~\citep{bhat2025certified}}

\noindent Third, we lift the mono-width restriction and obtain \textbf{Sound and complete solving for multi-width parametric bitvectors} (\cref{ch:multi-width}).
We prove that multi-width formulas over $n$ symbolic widths reduce, with linear blowup, to an equisatisfiable mono-width formula by encoding widths as bitmasks.
This translation collapses exponential enumeration into one fixed-width solver call and yields a new decidable fragment subsuming the prior state of the art.
\publishedat{Accepted with minor revision @ OOPSLA~2026~\citep{bhat2026multiwidth}}

\noindent Finally, we construct the first \textbf{Mechanized floating-point bitblaster} (\cref{ch:floating-point}).
We mechanize the SMT-LIB \qffp{} theory~\citep{smtlibfpa} and build verified bitblasting circuits for its operations
following SymFPU~\citep{symfpu}, and we evaluate the resulting solver on SMT-COMP~\citep{smtcomp2024} and on new
small-format (E4M3/E5M2) benchmarks~\citep{fp8fordl}.
\publishedat{Submitted @ POPL 2027~\citep{bhat2027floatingpoint}}

\paragraph{Publications and Attribution.}
Each contribution above is based on the publication whose venue is named beside it,
and the corresponding chapter reuses text from that paper.
All of these works are collaborative; throughout the thesis, ``we'' refers to myself and my coauthors,
and my own contributions are stated at the beginning of each chapter in an attribution note.
We also draw on ``Interactive Bitvector Reasoning using Verified Bit-Blasting''~\citep{leanbv} (OOPSLA 2025),
the fixed-width bitvector bitblasting paper on which I am second author,
and reuse text from it in the exposition of \bvdecide{} in the background chapter (\cref{ch:background}).
That work is presented as background rather than as a contribution of this thesis.

\paragraph{Plan of This Thesis.}
We begin by surveying the shared machinery of SAT solving, bitblasting, the Lean proof assistant,
and the verified \bvdecide{} bitblaster that the later procedures extend (\cref{ch:background}).
We then present the peephole-rewriting framework (\cref{ch:lean-mlir}).
With this framework in hand, we build solvers for parametric bitvectors,
first for a single symbolic width (\cref{ch:mono-width}), and then for many symbolic widths at once (\cref{ch:multi-width}).
We next turn to floating-point formulas and show how to bitblast them in a verified fashion,
yielding the first verified floating-point bitblaster (\cref{ch:floating-point}).
We conclude by charting unpublished results and future directions (\cref{ch:future-work}).
